\chapter{자주 묻는 질문과 함정}
\label{app:faq}

이 부록은 1권을 학습하며 독자가 자주 겪는 문제와 질문을 정리한 것이다. 각 항목은 실제 구현 중 발생할 수 있는 구체적 상황을 다룬다.

\section{환경 관련}

\subsection{Q: PyTorch 설치 후 CUDA가 인식되지 않아요}

\textbf{A}: 두 가지를 확인하라. 첫째, NVIDIA 드라이버가 설치되어 있는지 (\code{nvidia-smi} 실행). 둘째, PyTorch 버전과 CUDA 버전이 호환되는지. 예를 들어 CUDA 12.1용 PyTorch를 CUDA 11.8 시스템에 설치하면 인식 안 됨.

\begin{lstlisting}[style=bashstyle]
# CUDA 버전 확인
nvidia-smi | grep "CUDA Version"

# PyTorch 재설치 (CUDA 12.1용)
pip uninstall torch
pip install torch --index-url https://download.pytorch.org/whl/cu121
\end{lstlisting}

\subsection{Q: 학습이 너무 느려요}

\textbf{A}: 다음을 점검하라:
\begin{enumerate}
  \item GPU를 사용하고 있는지 (\code{device="cuda"} 설정).
  \item 배치 크기가 너무 작지 않은지 (GPU 활용도 낮음).
  \item 데이터 로딩이 병목인지 (\code{num\_workers} 증가).
  \item 혼합 정밀도(BF16)를 사용하고 있는지.
\end{enumerate}

\section{토크나이저 관련}

\subsection{Q: BPE로 학습한 토크나이저가 ``안녕''을 이상하게 토큰화해요}

\textbf{A}: 한국어가 코퍼스에 충분히 포함되지 않았기 때문. BPE는 빈도 기반이라 희귀 문자는 바이트로 분해된다. 해결책:
\begin{itemize}
  \item 한국어 코퍼스 비율을 높일 것.
  \item 어휘 크기를 늘릴 것 (한국어 음절 약 2만 개 + 서브워드).
  \item SentencePiece 같은 전문 토크나이저 사용.
\end{itemize}

\subsection{Q: encode/decode 라운드트립이 실패해요}

\textbf{A}: 공백 처리를 확인하라. BPE의 ``▁'' 문자가 디코딩 시 공백으로 변환되지 않을 수 있다. 우리의 구현에서는 \code{\_postprocess} 메서드가 이를 담당한다:

\begin{lstlisting}[language=Python]
def _postprocess(self, text):
    text = text.replace("▁", " ")
    return text.strip()
\end{lstlisting}

\subsection{Q: 특수 토큰 ID가 바뀌면 어떻게 되나요}

\textbf{A}: 치명적. 모델은 \code{<pad>=0, <bos>=1, <eos>=2, <unk>=3}을 가정하고 학습됐다. 토크나이저를 다시 학습하면서 특수 토큰 ID가 바뀌면, 모델 가중치가 더 이상 유효하지 않다. 항상 특수 토큰 ID를 고정하라.

\section{모델 관련}

\subsection{Q: 어텐션 출력이 NaN이 나와요}

\textbf{A}: 세 가지 원인이 흔하다:
\begin{enumerate}
  \item \textbf{소프트맥스 포화}: 로짓이 너무 커서 softmax가 0/1로 포화. 스케일링 ($\sqrt{d_k}$) 확인.
  \item \textbf{학습률 너무 큼}: 그래디언트 폭발. 학습률을 1/10으로 줄여보라.
  \item \textbf{FP16 오버플로우}: 활성값이 FP16 범위($\pm 65504$)를 초과. BF16으로 전환.
\end{enumerate}

\begin{lstlisting}[language=Python]
# NaN 확인
if torch.isnan(loss):
    print("NaN detected! Check learning rate and gradients.")
    print(f"Max grad: {max(p.grad.abs().max() for p in model.parameters() if p.grad is not None)}")
\end{lstlisting}

\subsection{Q: 모델이 같은 단어를 반복해요}

\textbf{A}: 샘플링 전략을 점검하라. Greedy는 반복 루프에 빠지기 쉽다. Top-P (P=0.9)로 전환하거나 Temperature를 0.7$\sim$0.9로 설정. 반복 패널티 (repetition penalty) 추가도 효과적.

\begin{lstlisting}[language=Python]
# 반복 패널티 (단순화)
def apply_repetition_penalty(logits, generated_ids, penalty=1.2):
    for token_id in set(generated_ids.tolist()):
        logits[token_id] /= penalty
    return logits
\end{lstlisting}

\subsection{Q: 컨텍스트 길이를 초과하면 에러가 나요}

\textbf{A}: 모델의 \code{context\_length}를 초과하는 입력을 넣으면 임베딩 단계에서 에러. \code{generate} 함수에서 슬라이딩 윈도우로 잘라내는 것을 확인하라:

\begin{lstlisting}[language=Python]
if input_ids.shape[1] > context_length:
    x = input_ids[:, -context_length:]  # 가장 최근 것만 유지
\end{lstlisting}

\section{학습 관련}

\subsection{Q: 손실이 줄어들지 않아요}

\textbf{A}: 다음을 점검하라:
\begin{enumerate}
  \item 학습률이 너무 작거나 (1e-6 이하) 너무 큼 (1e-2 이상).
  \item 데이터가 너무 적음 (최소 수천 토큰 필요).
  \item 옵티마이저가 잘못됨 (AdamW의 beta, eps 확인).
  \item 그래디언트가 흐르지 않음 (requires\_grad 확인).
\end{enumerate}

\begin{lstlisting}[language=Python]
# 그래디언트 흐름 확인
for name, p in model.named_parameters():
    if p.requires_grad and p.grad is None:
        print(f"No gradient for {name}")
    elif p.grad is not None and p.grad.abs().max() < 1e-10:
        print(f"Vanishing gradient for {name}")
\end{lstlisting}

\subsection{Q: 학습 초기에 손실이 발산해요}

\textbf{A}: Warmup이 부족하거나 초기화가 잘못됐을 가능성. 다음을 시도:
\begin{itemize}
  \item Warmup 스텝 늘리기 (100$\sim$1000).
  \item 그래디언트 클리핑 확인 (\code{grad\_clip=1.0}).
  \item 초기화 표준편차 줄이기 (0.02 $\to$ 0.01).
\end{itemize}

\subsection{Q: 체크포인트를 로드하면 다른 결과가 나와요}

\textbf{A}: 두 가지 원인:
\begin{enumerate}
  \item \textbf{시드 미설정}: \code{set\_seed(42)}를 호출했는지 확인.
  \item \textbf{드롭아웃}: 모델을 \code{eval()} 모드로 전환했는지 확인.
  \item \textbf{CuDNN 비결정론}: \code{torch.backends.cudnn.deterministic = True}.
\end{enumerate}

\section{분산 학습 관련 (2권 미리보기)}

\subsection{Q: DDP로 학습하면 속도가 오히려 느려요}

\textbf{A}: 통신 병목일 가능성. 다음을 점검:
\begin{itemize}
  \item 배치 크기가 너무 작아 연산/통신 비율이 낮음.
  \item NVLink가 아닌 느린 통신 (PCIe 등) 사용.
  \item 그래디언트 동기화가 매 스텝 일어남 (gradient accumulation으로 줄이기).
\end{itemize}

\subsection{Q: FSDP에서 OOM이 나요}

\textbf{A}: 활성화 메모리가 부족. 활성화 체크포인팅을 켜거나 배치 크기를 줄이기. CPU offload도 고려.

\section{디버깅 팁}

\subsection{단계별로 검증하라}

복잡한 버그는 작은 단위로 쪼개서 검증. 토크나이저만 단독 테스트, 어텐션만 단독 테스트 등.

\begin{lstlisting}[language=Python]
# 토크나이저 단독 테스트
tok = BPETokenizer()
tok.train(["hello world"], vocab_size=50)
print(tok.encode("hello"))
print(tok.decode(tok.encode("hello")))

# 어텐션 단독 테스트
attn = MultiHeadAttention(config)
x = torch.randn(1, 4, 32)
out = attn(x)
print(out.shape, out.dtype, torch.isnan(out).any())
\end{lstlisting}

\subsection{shape를 항상 출력하라}

버그의 80\%는 shape 불일치. 모든 단계에서 shape를 출력하는 습관.

\begin{lstlisting}[language=Python]
print(f"input shape: {x.shape}")
print(f"after embedding: {emb.shape}")
print(f"after attention: {attn_out.shape}")
\end{lstlisting}

\subsection{작은 입력으로 시작하라}

배치 1, 시퀀스 4로 시작해서 동작을 확인한 후 크기를 키울 것. 큰 입력에서 디버깅하면 출력이 너무 많아 혼란.

\section{성능 최적화 팁}

\subsection{혼합 정밀도 사용}

\begin{lstlisting}[language=Python]
from torch.cuda.amp import autocast, GradScaler

scaler = GradScaler()
with autocast(dtype=torch.bfloat16):
    logits = model(inputs)
    loss = criterion(logits, targets)
scaler.scale(loss).backward()
scaler.step(optimizer)
scaler.update()
\end{lstlisting}

BF16은 FP16과 달리 scaler가 필요 없다. A100 이상에서 권장.

\subsection{torch.compile}

PyTorch 2.0의 \code{torch.compile}로 모델을 컴파일하면 1.5$\sim$2배 속도 향상.

\begin{lstlisting}[language=Python]
model = GPT(config)
model = torch.compile(model)  # 컴파일
\end{lstlisting}

\subsection{프로파일링}

\begin{lstlisting}[language=Python]
with torch.profiler.profile(activities=[torch.profiler.ProfilerActivity.CPU,
                                         torch.profiler.ProfilerActivity.CUDA]) as prof:
    # 학습 코드
prof.export_chrome_trace("trace.json")
\end{lstlisting}

Chrome 브라우저에서 \code{chrome://tracing}으로 열어 병목 분석.
